\documentclass[12pt,a4paper]{article}
\usepackage[T2A]{fontenc}
\usepackage[utf8]{inputenc}
\usepackage[russian]{babel}
\usepackage{geometry}
\usepackage{amsmath}
\usepackage{enumitem}
\usepackage{listings}
\geometry{left=23mm,right=20mm,top=19mm,bottom=20mm}
\setlength{\parindent}{1.2cm}
\setlength{\parskip}{0.35em}
\lstset{basicstyle=\ttfamily\small,frame=single,breaklines=true,showstringspaces=false}
\begin{document}
\begin{center}
{\Large\bfseries Отчёт по программированию}\\[4mm]
{\Large\bfseries Вариант 38. Служба доставки беспилотниками}\\[6mm]
ФИО: Орлов Фёдор Алексеевич\\[3mm]
Группа: БПИ253
\end{center}

\section*{1. Постановка задачи}
Разработано консольное приложение на языке C для моделирования службы доставки.
Есть склад, заказы с массой и сроками, беспилотники с ограничениями по грузоподъёмности и энергии, зарядные станции с ограниченным числом мест, запрещённые клетки карты и погода. Диспетчер распределяет заказы между беспилотниками. Программа показывает происходящие события и статистику. Для моделирования используется библиотека \texttt{pthread}.

\section*{2. Сценарий и взаимодействие потоков}
Моделирование разбито на дискретные такты. За каждый такт выполняются четыре этапа:
\begin{enumerate}[nosep]
\item Потоки заказов проверяют время поступления и сроки.
\item Поток диспетчера назначает свободным беспилотникам подходящие заказы.
\item Потоки беспилотников перемещают аппараты на один шаг, получают груз, доставляют его или возвращаются на зарядку.
\item Потоки станций пополняют энергию заряжающихся аппаратов.
\end{enumerate}
Для этапа создаются потоки функцией \texttt{pthread\_create}, затем основной поток вызывает \texttt{pthread\_join}. Для защиты общих структур применяется один \texttt{pthread\_mutex\_t}.

\section*{3. Основные алгоритмы}
\textbf{Построение пути.} Карта представлена двумерной сеткой. Путь между двумя точками ищется обходом в ширину (BFS): посещаем соседние доступные клетки и запоминаем предка каждой клетки. Если путь найден, восстанавливаем его от назначения к началу. Для карты из $R$ строк и $C$ столбцов время и память алгоритма: $O(RC)$.

\textbf{Назначение заказа.} Диспетчер перебирает свободные аппараты. Аппарат должен поднять заданную массу, успеть до срока при хорошей погоде и иметь запас энергии на три маршрута: от текущей позиции к складу, от склада к клиенту, от клиента к ближайшей доступной по карте зарядной станции. Выбирается аппарат с минимальным расчётным временем доставки. Если подходящего аппарата пока нет, заказ остаётся в ожидании до следующего такта или истечения срока.

\textbf{Зарядка.} После доставки беспилотник направляется к достижимой станции. Если все места заняты, ждёт на станции и повторяет попытку на следующем такте. Скорость зарядки фиксирована. Зарядка выполняется и при плохой погоде.

\textbf{Погода.} В модели есть хорошая и плохая погода. В плохую погоду полёты приостанавливаются. Когда погода улучшается, маршрут пересчитывается BFS, полёт продолжается.

\section*{4. Интерфейс и запуск}
Приложение компилируется командой \texttt{make}. Запуск без аргументов открывает меню: старт моделирования, печать карты, просмотр беспилотников/станций/заказов, выход. Для автоматического запуска используется:
\begin{lstlisting}
./fly_simulator --auto
./fly_simulator --auto --scenario scenarios/small.txt
./fly_simulator --auto --delay 100
make test
\end{lstlisting}
Входные данные читаются из текстовых сценариев. Формат описан в \texttt{Readme.md}. Указание \texttt{max\_time 0} позволяет моделировать без заранее заданного лимита времени; Ctrl+C прерывает выполнение.

\section*{5. Вывод и результаты проверки}
Программа выводит получение и назначение заказа, взятие груза, движение и изменение энергии, смену погоды, пересчёт пути, прибытие, зарядку, возврат и просрочку. В конце выводятся количество доставленных и неисполненных заказов, пройденное расстояние и доля успешных доставок.

На сценарии \texttt{scenarios/demo.txt} получилось: 6 заказов, 4 доставлены, 2 завершились отказом, суммарный путь --- 163 клетки. Проверяются ситуации с препятствиями, нехваткой энергии, вместимостью станции, погодными паузами и истечением срока. Автоматических тестов --- 10; запускаются командой \texttt{make test}.

\section*{6. Вывод}
Получено простое консольное многопоточное приложение, соответствующее основным требованиям варианта 38. Для учебной наглядности оставлены один способ маршрутизации (BFS), одна стратегия назначения и один общий мьютекс. 
\end{document}
